% engines: lualatex, pdflatex
% expect-warnings: 0
% expect-errors: 0
\documentclass{article}
\usepackage[margin=1cm]{geometry}
\usepackage{../../latex/pgf-go}
\newcommand\case[1]{\par\medskip\noindent\textbf{#1}\par\nopagebreak}
\pagestyle{empty}
\begin{document}

% draw = region with polygons, so it works with both kernels.

\case{A hole (D4-F6 are white stones) / two territories that touch at a corner (D4) / two \texttt{\textbackslash territory} of the same color are one region}
\begin{goban}[label = {at = {above, left}}]
	\stones[w]{D4:F6}
	\territory[b, draw = region]{B2,B8,H8,H2}
\end{goban}
\begin{goban}[label = {at = {above, left}}]
	\territory[b, draw = region]{A1,C1,C3,A3}
	\territory[b, draw = region]{E5}
\end{goban}
\begin{goban}[label = {at = {above, left}}]
	\territory[w, draw = region]{A9,D9,D6,A6}
	\territory[w, draw = region]{E9,E6}
	\territory[b, draw = region]{A1,I1,I3,A3}
\end{goban}

\pgfgoregion{b}{polygons} polygons (1), \pgfgoregion{w}{polygons} (1),
area \pgfgoregion{w}{area} (20), perimeter \pgfgoregion{w}{perimeter} (18),
vertices of polygon 1: \pgfgoregion{w}{polygon 1} (A5,F5,F9,A9).

\case{An island in a hole (13x13; C3-K11 and D4-J10 are white): line and vertices}
\begin{goban}[13, label = {at = {above, left}}]
	\stones[w]{C3 r K11, D4 r J10}
	\territory[b, draw = region, line, vertices]{A1,A13,N13,N1}
\end{goban}

polygons \pgfgoregion{b}{polygons} (3), outlines \pgfgoregion{black}{outlines} (2),
holes \pgfgoregion{b}{holes} (1), vertices \pgfgoregion{b}{vertices} (12),
area \pgfgoregion{b}{area} (116), perimeter \pgfgoregion{b}{perimeter} (104),
polygon 2: \pgfgoregion{b}{polygon 2} (C3,C11,K11,K3), polygon 4: [\pgfgoregion{b}{polygon 4}] (empty),
white: \pgfgoregion{w}{polygons} (0), neutral: \pgfgoregion{n}{area} (0).

\case{Line styles: below the stones / red, without fill / vertices without line; the last style of each color wins}
\begin{goban}
	\stones[b]{C3:G7}
	\territory[b, draw = region, line = {true, above = false, width = 3pt}]{B2,B8,H8,H2}
\end{goban}
\begin{goban}
	\stones[b]{D4:F6}
	\territory[b, draw = region, fill = {opacity = 0}, line = {true, color = red}]{B2,B8,H8,H2}
\end{goban}
\begin{goban}
	\territory[w, draw = region, line]{A1,E1,E5,A5}
	\territory[w, draw = region, vertices = {true, radius = 3pt, color = blue}]{H8}
\end{goban}

\case{Partial board (A1-E5): the region is clipped, the line ends like the grid / marks on the vertices}
\begin{goban}[partial = A1:E5, label = {at = {above, left}}]
	\territory[b, draw = region, line, vertices]{A1,G1,G3,C3,C7,A7}
\end{goban}
\begin{goban}
	\territory[b, draw = region]{A1,G1,G3,C3,C7,A7}
	\marks[triangle]{\pgfgoregion{b}{polygon 1}}
\end{goban}

\case{Players style; neutral is filled point by point; name and resume}
{\pgfgoset{players = {black = {territory = {fill = {color = blue}, line}}}, territory = {draw = region}}%
\begin{goban}[5, name = region]
	\stones[b]{B1:5}\stones[w]{D1:5}
	\territory[b]{A1,A5}
	\territory[n]{C1,C5}
	\territory[w, forget]{E1,E5}
\end{goban}
\begin{goban}[5, resume = {region, territory}]\end{goban}}

\end{document}
